\section{结论}\label{sec:conclusion}

本文综述了挂谷猜想从1917年转针问题到2025年三维证明的百年历程。可以概括出三条主线：

\begin{enumerate}[label=(\arabic*)]
    \item \textbf{几何与分析主线}：从 Besicovitch 的 Perron 树（面积归零），到 Davies 的平面满维数，再到 Córdoba 的 $L^2$ 计数与 Wolff 的毛刷论证——挂谷集的研究把“测度为零的病态对象”锻造成调和分析中波包交叠的标准模型。
    \item \textbf{组合与代数主线}：Dvir 的有限域多项式方法排除了纯组合障碍并辐射到整个离散数学；Guth 的多项式分划则把代数几何工具带回实数问题，补齐多线性端点。两条支流最终汇聚于 Wang--Zahl 对三维的完整证明——管族的结构二分（粘性/颗粒）是理解“尺度递归”的成熟范式。
    \item \textbf{网络效应}：挂谷猜想位于局部光滑化、限制性、Bochner--Riesz 一整族猜想的组合核心；三维的解决立刻打通了三维情形的整条链条，而 $n \ge 4$ 的挂谷问题仍是这些领域共同的试金石。
\end{enumerate}

挂谷猜想的故事再次印证了数学的深刻统一性：一个关于“旋转一根针”的初等谜题，最终答案藏在波包的相消干涉与代数几何的分划结构之中。悬而未决的高维情形，正是这一统一性视野下的下一个战场。
